%=====================================================================
% Datos · Estados y categorías
%   La tabla de cobertura la genera bd/exportar_datos.py a partir de las
%   restricciones CHECK de la base.
%=====================================================================
\subsection*{Estados y categorías}
\addcontentsline{toc}{subsection}{Estados y categorías}

Casi todas las entidades del modelo tienen un estado, un tipo o una categoría que decide qué columnas aplican y qué puede pasarle a la fila. La captura 4 cuenta las filas de cada valor en las dieciocho columnas más importantes: las cuentas están en los cinco estados y tienen los tres roles, las partidas son de los tres tipos y terminan de las cinco formas, y los reportes, las sanciones, las cajas y las invitaciones pasan por todos sus estados.

\captura{04-estados}{Captura 4. Filas por valor en las columnas de estado y de categoría.}

La tabla siguiente hace la misma cuenta para todas las columnas de dominio cerrado: las \dnNumDominios{} cuyos valores fija una restricción \at{CHECK}. \at{bd/exportar\_datos.py} lee esos valores del catálogo del sistema, así que la tabla no puede omitir ninguno. \dnNumDominiosCompletos{} columnas tienen registros en todos sus valores, y en total \dnNumValoresConFilas{} de los \dnNumValores{} valores tienen al menos uno.

% Archivo generado por bd/exportar_datos.py. No editar a mano.
\begin{center}\footnotesize\setlength{\tabcolsep}{2pt}
\begin{longtable}{|L{143.2pt}|L{159.3pt}|L{155.2pt}|}
\hline
\textbf{\hspace{0pt}Columna} & \textbf{\hspace{0pt}Valores con filas (filas)} & \textbf{\hspace{0pt}Valores sin filas} \\ \hline
\endfirsthead
\hline
\textbf{\hspace{0pt}Columna} & \textbf{\hspace{0pt}Valores con filas (filas)} & \textbf{\hspace{0pt}Valores sin filas} \\ \hline
\endhead
ObjetoCosmetico.rareza & LEGENDARIO (2), EPICO (5), RARO (7), COMUN (9) & — \\ \hline
PalabraProhibida.ambito & AMBOS (3), CHAT (2), NICKNAME (3) & — \\ \hline
PalabraProhibida.idioma & en (3), es-MX (4), NULL (1) & — \\ \hline
Usuario.estado\_\allowbreak{}cuenta & ELIMINADA (1), BANEADA (1), SUSPENDIDA (1), ACTIVA (7), PENDIENTE (1) & — \\ \hline
Usuario.idioma\_\allowbreak{}preferido & en (1), es-MX (10) & — \\ \hline
Usuario.rol & ADMINISTRADOR (1), MODERADOR (1), JUGADOR (9) & — \\ \hline
Usuario.tipo\_\allowbreak{}cuenta & REGISTRADA (10), INVITADO (1) & — \\ \hline
Sesion.motivo\_\allowbreak{}cierre & SANCION (2), BAJA\_\allowbreak{}CUENTA (1), CAMBIO\_\allowbreak{}CREDENCIALES (1), CIERRE\_\allowbreak{}REMOTO (1), CIERRE\_\allowbreak{}VOLUNTARIO (2), NULL (8) & — \\ \hline
CodigoSegundoFactor.estado & INVALIDADO (1), CONSUMIDO (2), PENDIENTE (1) & — \\ \hline
TokenVerificacionCorreo.estado & INVALIDADO (1), USADO (9), PENDIENTE (2) & — \\ \hline
TokenVerificacionCorreo.proposito & ALTA (11), CAMBIO\_\allowbreak{}CORREO (1) & — \\ \hline
TokenRecuperacion.estado & INVALIDADO (1), USADO (1), PENDIENTE (1) & — \\ \hline
AceptacionTerminos.idioma & en (1), es-MX (10) & — \\ \hline
Avatar.formato & PNG (1), JPG (2) & — \\ \hline
UsuarioObjeto.origen & NIVEL (1), CAJA (2), COMPRA (4), INICIAL (88) & — \\ \hline
Partida.estado & FINALIZADA (11), EN\_\allowbreak{}CURSO (1) & PENDIENTE\_\allowbreak{}DE\_\allowbreak{}CIERRE \\ \hline
Partida.forma\_\allowbreak{}termino & ABANDONO (1), TABLAS (1), RENDICION (2), TIEMPO\_\allowbreak{}AGOTADO (1), META (6), NULL (1) & — \\ \hline
Partida.tipo & IA (1), PRIVADA (1), CLASIFICATORIA (10) & — \\ \hline
Participacion.forma\_\allowbreak{}termino & ABANDONO (3), NULL (22) & RENDICION, TIEMPO\_\allowbreak{}AGOTADO, META \\ \hline
Participacion.resultado & TABLAS (2), PERDIDA (12), GANADA (9), NULL (2) & — \\ \hline
Jugada.orientacion & VERTICAL (6), HORIZONTAL (17), NULL (111) & — \\ \hline
Jugada.tipo & MOVIMIENTO (111), MURO (23) & — \\ \hline
OfertaTablas.estado & CADUCADA (1), RECHAZADA (1), ACEPTADA (1), PENDIENTE (1) & — \\ \hline
Sala.estado & CERRADA (1), ABIERTA (1) & EN\_\allowbreak{}PARTIDA \\ \hline
Invitacion.estado & EXPIRADA (1), RECHAZADA (1), ACEPTADA (1), PENDIENTE (1) & — \\ \hline
Mensaje.canal & ESPECTADORES (1), PARTIDA (4), SALA (3), AMIGOS (1), GLOBAL (2) & — \\ \hline
Caja.estado & CADUCADA (1), ABIERTA (1), SIN\_\allowbreak{}ABRIR (1) & — \\ \hline
Caja.origen & PARTIDA (2), COMPRA (1) & — \\ \hline
MovimientoMoneda.tipo & CONVERSION (1), COMPRA\_\allowbreak{}CAJA (1), COMPRA (4), NIVEL (4), PARTIDA (20) & DEVOLUCION \\ \hline
Reporte.estado & PENDIENTE (1), EN\_\allowbreak{}REVISION (1), RESUELTO\_\allowbreak{}CON\_\allowbreak{}SANCION (2), RESUELTO\_\allowbreak{}SIN\_\allowbreak{}SANCION (1) & — \\ \hline
Sancion.ambito & CHAT (2), CUENTA (3) & — \\ \hline
Sancion.tipo & TEMPORAL (3), PERMANENTE (2) & — \\ \hline
Apelacion.estado & PENDIENTE (1), RECHAZADA (1), ACEPTADA (1) & EN\_\allowbreak{}REVISION \\ \hline
BitacoraModeracion.accion & INTENTO\_\allowbreak{}SIN\_\allowbreak{}PERMISO (1), CONSULTA\_\allowbreak{}BITACORA (1), ROL\_\allowbreak{}OTORGADO (1), APELACION\_\allowbreak{}RESUELTA (2), SANCION\_\allowbreak{}APLICADA (5), REPORTE\_\allowbreak{}RESUELTO (3), REPORTE\_\allowbreak{}LIBERADO (1), REPORTE\_\allowbreak{}TOMADO (5) & ROL\_\allowbreak{}RETIRADO \\ \hline
BitacoraAcceso.resultado & CUENTA\_\allowbreak{}ELIMINADA (1), SEGUNDO\_\allowbreak{}FACTOR\_\allowbreak{}ACTIVADO (1), CAMBIO\_\allowbreak{}CORREO\_\allowbreak{}SOLICITADO (1), CAMBIO\_\allowbreak{}CONTRASENA (1), CIERRE\_\allowbreak{}REMOTO (1), CIERRE\_\allowbreak{}VOLUNTARIO (2), RECUPERACION\_\allowbreak{}EXITOSA (1), RECUPERACION\_\allowbreak{}CORREO\_\allowbreak{}INEXISTENTE (1), RECUPERACION\_\allowbreak{}SOLICITADA (3), CUENTA\_\allowbreak{}VERIFICADA (9), ALTA\_\allowbreak{}INVITADO (1), REGISTRO\_\allowbreak{}RECHAZADO (1), REGISTRO\_\allowbreak{}EXITOSO (10), SEGUNDO\_\allowbreak{}FACTOR\_\allowbreak{}FALLIDO (1), CUENTA\_\allowbreak{}SANCIONADA (2), CUENTA\_\allowbreak{}PENDIENTE (1), CONTRASENA\_\allowbreak{}INCORRECTA (3), BLOQUEO\_\allowbreak{}VIGENTE (1), USUARIO\_\allowbreak{}INEXISTENTE (1), EXITO (12) & SEGUNDO\_\allowbreak{}FACTOR\_\allowbreak{}DESACTIVADO, CAMBIO\_\allowbreak{}CORREO\_\allowbreak{}CONFIRMADO, VINCULACION\_\allowbreak{}EXITOSA, REGISTRO\_\allowbreak{}SIN\_\allowbreak{}CORREO, SUSPENSION\_\allowbreak{}VENCIDA, ABANDONADO \\ \hline
\end{longtable}
\end{center}


Los valores sin registros son estados de excepción o de paso, o eventos que no ocurren en la fotografía; incluirlos obligaría a romper la coherencia del escenario:

\begin{reglas}
  \item \textbf{\at{PENDIENTE\_DE\_CIERRE} de \textit{Partida}.} Una partida queda así solo si se agotan los reintentos al cerrarla (\mbox{CU-25} \mbox{EX-01}), y el servidor la cierra en cuanto puede.
  \item \textbf{\at{EN\_PARTIDA} de \textit{Sala}.} Una sala lo está solo mientras dura su partida privada (\mbox{CU-18}). La de la sala 1 ya terminó, y la única partida en curso es clasificatoria.
  \item \textbf{\at{META}, \at{TIEMPO\_AGOTADO} y \at{RENDICION} en la salida de una \textit{Participacion}.} Esa columna solo se escribe cuando un jugador sale antes que los demás de una partida de cuatro, y en la de prueba los tres salen por abandono. Las cinco formas sí aparecen como forma de término de las partidas.
  \item \textbf{\at{EN\_REVISION} de \textit{Apelacion}.} Una apelación tomada vuelve a la cola a los treinta minutos (\mbox{CU-43} \mbox{RN-03}), y a la hora de la fotografía la única moderadora conectada es la que aplicó las sanciones apeladas, que no puede resolverlas (\mbox{CU-45} \mbox{RN-06}).
  \item \textbf{\at{DEVOLUCION} de \textit{MovimientoMoneda}.} Solo aparece si falla la confirmación de la compra de una caja (\mbox{CU-39} \mbox{EX-02}) o para corregir un movimiento (\mbox{CU-40} \mbox{RN-01}), y en el escenario no falla ninguna compra.
  \item \textbf{Acciones de las bitácoras.} Nadie retira un rol, desactiva el segundo factor, confirma un cambio de correo, vincula su cuenta de invitado ni abandona un inicio de sesión a medias; el alta no encuentra caído el servicio de correo, y la única suspensión vigente aún no vence.
\end{reglas}
